\chapter{Ausblick}

Die nächsten möglichen Schritte für das Projekt lässt sich in unterschiedliche Kategorien unterteilen. Die Priorität
könnte dabei auf \enquote{Iteratives Deeping} und anschliessend \enquote{UCI Support} gelegt werden. \\
Dadurch können dann Testing Frameworks wie \href{https://github.com/AndyGrant/OpenBench}{OpenBench} eingesetzt werden,
welche ein UCI Support vorbedingen. Solche Testing Frameworks verwenden
\href{https://www.chessprogramming.org/Sequential_Probability_Ratio_Test}{Sequential Probability Ratio Tests}
die anhand von ganzen Spielen Weiterentwicklungen statistisch überprüfen (\cite{sequentialProbabilityRatioTest}) \\
Zusätzlich kann mittels Iteratives Deeping und UCI Support standardisiert Spiele gegen andere Bots durchgeführt werden
und so an Turniere teilgenommen werden.

\section{Zuqualität}

Wie es sich bei der Evaluation der Zugqualität herausgestellt hat, hat die Engine Verbesserungspotenzial im \enquote{positionale Stellungen}.
Hierfür können Heuristiken wie \href{https://www.chessprogramming.org/Pawn_Structure}{Pawn Structure}
\href{https://www.chessprogramming.org/Mobility}{Mobility} oder \href{https://www.chessprogramming.org/King_Safety}{King Safety}
implementiert werden.

\section{Suchalgorithmus}



\section{Weitere Funktionalitäten}

\subsection{Endgame Tablebase \& Opening Books}

- Endgames https://www.chessprogramming.org/Endgame_Tablebases
- Openning Books https://www.chessprogramming.org/Opening_Book

\subsection{Iterative Deepening}

https://www.chessprogramming.org/Iterative_Deepening

\subsection{UCI}

\subsection{}
